\documentclass[stu,12pt,floatsintext]{apa7}

\usepackage[american]{babel}
\usepackage{booktabs}

\usepackage{csquotes} 
\usepackage[style=apa,sortcites=true,sorting=nyt,backend=biber]{biblatex}
\DeclareLanguageMapping{american}{american-apa} 
\addbibresource{bibliography.bib}

\usepackage[T1]{fontenc} 
\usepackage{mathptmx} 
\usepackage{hyperref}
\title{Hyperparameter Optimization in Artificial Intelligence}
\author{Paul Beggs}
\duedate{December, 15, 2025}
\affiliation{Hendrix College}
\course{CSCI 335: Artificial Intelligence}
\professor{Dr. Gabriel Ferrer} 

\begin{document}
\maketitle

\renewcommand{\labelenumi}{\theenumi.}

In a nutshell, we use computers to solve problems through programs in computer science. However, it is common to encounter situations where we cannot write a traditional program to solve a challenging problem. These problems can have a search space that is too large (e.g., checkers or chess), or must be generalized so that we can apply a solution to cases that the program has not seen before. Traditionally, these problems are best suited to artificial intelligence, as these algorithms can handle them effectively using sophisticated methods. Some of these algorithms employ hyperparameters, or values set before training \parencite{weerts2020importance}. In this paper, we show that as the complexity of training an algorithm increases, the need for more complex hyperparameter optimization methods also increases.

We first dive deeper into hyperparameters and how they are used in machine learning so that artificial intelligence algorithms work as intended. As mentioned before, a hyperparameter is a value set by a researcher and used during training. A typical example is tuning the \(k\) parameter in the k-nearest neighbors algorithm. Depending on this value, we can set the threshold for grouping data points with similar values, but more on this later. Tuning hyperparameters is commonly done by trial-and-error: a researcher would provide a good guess of the optimal hyperparameter, then train the model using machine learning, and review the results to see whether the change affected the model's predictions \parencite{bergstra2012random}. However, this manual process can be tedious and frustrating, as finding the best value essentially requires a good guess. This problem is exacerbated when the algorithm's complexity increases due to more hyperparameters or longer learning time, thereby adding additional time before the researcher can try another random pair of hyperparameters.

In the domain of artificial intelligence, there is a phenomenon known as ``the curse of dimensionality.'' In which, as we add more data that we want to track, the harder it is to meaningfully identify relationships \parencite{geeksforgeeks_2024_curse}. This problem can sometimes be addressed by employing dimensionality reduction techniques, but we will not explore that in this paper; instead, we assume that all hyperparameters are required. Thus, for algorithms with many hyperparameters, such as neural networks, we must use more sophisticated methods to optimize them. 

Thus, we introduce four hyperparameter optimizers: grid search, random search, Bayesian optimization, and simulated annealing. We will then investigate three distinct problems in artificial intelligence: handwriting classification, sentiment analysis, and image classification, and how each optimizer can be applied to the problem. Depending on the application, we will determine if one or more optimizers are best suited for the situation, and walk through why that is the case. 

\section{Hyperparameter Optimizers}

\subsection{Grid Search}

One of the most common and straightforward hyperparameter optimizers is grid search. This optimizer works by performing an exhaustive search over a domain and trying every possible combination of hyperparameters. The search starts by having a researcher manually define points that are equally spaced across the domain. Then, they run a loop for each hyperparameter choice and compare the predicted values to the actual values for each combination. The hyperparameters that yield the best accuracy are selected for the model. Note that there are many metrics for measuring accuracy, such as root mean square error or negative log-likelihood. However, we will not cover them in this paper and instead assume a generic accuracy metric.

Since hyperparameters can take the form of any variable, they range from having infinite possibilities (numerical) to single-digit possibilities (categorical). For example, some continuous variables could be \(c\) and \(\gamma\) (regularization hyperparameters), or linear and RBF (kernel type) for support vector machines \parencite{scikitlearn_2019_sklearnsvmsvc}. On the one hand, we would have to search over an arbitrarily large number of values of \(c\) and \(\gamma\), but on the other hand, we only need to search over five kernels. Hence, categorical variables are where grid search shows the most promise with exhaustive search. 

From the grid search's description, it is not hard to imagine that as the number of hyperparameters increases, the search time in the space rises as well. This increase is due to the exponential time complexity of running loops for every hyperparameter. That is, if we have \(k\) hyperparameters, and we decide to test \(n\) distinct values, the total number of evaluations is \(n^{k}\). To put that into practice, consider the following example: we have four parameters with 10 values, which directly translates to \(10^{4} =\) 10,000 training runs. If each run takes 5 minutes, the search would take over a month. 

While running 10,000 training runs can be a significant roadblock for most applications, grid search has a highly beneficial property of being ``embarrassingly parallel.'' Because the results of each set of hyperparameters are independent of each other (i.e., set one's results do not affect set two’s), we can run multiple hyperparameter configurations at the same time across various systems \parencite{bergstra2012random}. Essentially, if we have 10 GPUs, we can train 10 models simultaneously, reducing the wait time by 90\%. This makes grid search viable for organizations that have massive computational resources, even when the time complexity is exponential. However, other organizations may prefer an option that is just as good as grid search but does not suffer from the curse of dimensionality, thanks to its low effective dimensionality. That option is random search. 

\subsection{Random Search}

Random search, like grid search, has a researcher choose candidate hyperparameter values from a grid to maximize a model's predictive performance. However, it differs in how these points are determined and does not perform exhaustive searches over the entire domain. Instead, random search picks hyperparameters based on a uniform distribution \parencite{pilario2020kernel}. Given this random sampling of parameters, if we use the same number of points as grid search, random search finds the hyperparameters that yield the highest success rate for the model \parencite{bergstra2012random}.

These performance gains stem from the fact that random search does not use a fixed grid structure. We initialize the search by picking candidate points from our model's hyperparameter sets. Still, we utilize a uniform distribution to give every possible hyperparameter within a preset range an equally likely probability of being picked \parencite{zabinsky2009random}. If our success metric reports higher values for a particular set of hyperparameters, the algorithm reports them, but does not use that information for later candidates. We continue the search process until we reach a cutoff threshold for the success metric, exhaust the candidate points, or reach a budget cap.

To show the advantage of random search over grid search, consider a scenario where we can utilize ``low effective dimensionality'' \parencite{bergstra2012random}. While tuning a neural network, we determine that only two hyperparameters are of importance: learning rate and momentum. The learning rate controls the size of each step taken toward a minimum, whereas momentum influences how quickly the gradient accelerates in that direction \parencite{orr_1999_momentum}. In this example, we will assume that the learning rate is the more critical parameter over momentum. If we used a \(3 \times 3\) grid search, we would run 9 tests, but because we reused values we had already tested, we effectively used only 3 distinct values for the learning rate. We wasted six trials on those same learning rate values, with momentum values that were less important. If we instead used random search for the same nine trials, we would try nine distinct values for the learning rate, thereby casting a wider net for both hyperparameters. This example shows that random search is much more efficient in high-dimensional spaces as it samples important hyperparameters more effectively.

We have explored how random search can be better than grid search, but that does not mean random search is without its limitations. Since it, like grid search, treats sampled hyperparameters as independent, it does not ``learn'' from its past exploration. That is, if one hyperparameter yields a great accuracy score, it does not focus on that region. Instead, it will randomly sample another hyperparameter. The lack of learning from experience allows for the development of more sophisticated methods, such as Bayesian optimization, which uses experience as a fundamental part of the algorithm.  

\subsection{Bayesian Optimization}

Unlike grid search and random search, which do not take previous data into account, the Bayesian optimization process learns from previous data and incorporates them into a probabilistic model of an ``objective function''\footnote{We can also think of the objective function as our generic success metric that we have been using.} \parencite{frazier2018tutorial}. The objective function is generally considered to be continuous, but it does not have linearity, concavity, and ``first- or second-order derivatives'' \parencite{frazier2018tutorial}. These qualities make the objective function into a ``black box problem,'' where we essentially do not know what its qualities are, and where every evaluation is computationally expensive \parencite{WU201926}. So, we are trying to predict hyperparameters using a probabilistic model for a function that we know very little about. Thus, with the Bayesian approach, we can make more informed guesses that perform better than other optimization methods \parencite{WU201926}.

Predicting the objective function is similar to how Bayesian statisticians model posterior distributions. They start with a ``hunch'' of what they think the distribution may look like, called the prior distribution. This hunch comes from previous experiences with the objective function, and does not have to be based on fact \parencite{WU201926}. Instead, a researcher can use their opinion for what they think the prior distribution would look like.  Then they introduce new data to the model and calculate a new posterior distribution, which provides information on where to sample next. Through this process of educated guessing, we can approximate the optimal point of the objective function, corresponding to the optimal set of hyperparameters that maximizes it.

After the posterior distribution is updated, the algorithm uses an acquisition function to determine which hyperparameter configuration to test next \parencite{WU201926}. There are many acquisition functions available, but the one we will discuss here is the upper-confidence bound (UCB) acquisition function, which balances exploration and exploitation. On one hand, exploration is the act of exploring areas of high variance, where the function could be much higher or lower. By sampling these regions of high variance, we are exploring possible points that are better than the current point's neighbors. Exploitation, on the other hand, only focuses on finding high probabilities of success, or the most significant point on the graph. With UCB, the balance between exploration and exploitation is controlled using a hyperparameter \(k\). The higher \(k\) is, the more it values exploration; the lower it is, the more it values exploitation.

\subsection{Simulated Annealing}

The simulated annealing algorithm is based on the physical process of material annealing. That is, heating a material to a high temperature, then cooling the material within a tightly controlled process that allows the atoms to be arranged in an optimal pattern \parencite{rutenbar2002simulated}. For hyperparameter optimization, we start by exploring the space by trying any hyperparameter configuration, regardless of improvements for our success metric. Then the algorithm becomes more selective in the hyperparameters it chooses, eventually settling on the configuration it deems best. 

One of the first simulated annealing algorithms developed is the ``Metropolis algorithm,'' which dictates how the system explores and eventually settles \parencite{rutenbar2002simulated}. It works by initializing a temperature \(T\) and a stopping criterion, \(M\), for which \(T\) slowly decrements, and a counter enumerates until it reaches \(M\). For each loop, we randomly pick a hyperparameter configuration close to the current position, and check the change in the system's ``energy,'' denoted as \(\Delta E\). This energy value is related to the model's error rate, which should be minimized.

If \(\Delta E\) is negative (the new error is lower), we approve of the change and move toward a minimum. However, if \(\Delta E\) is positive (the new error is higher), the algorithm does not immediately reject the move. Instead, it accepts this upward move based on the Boltzmann probability \(P = e^{\Delta E/T}\). From this equation, we can see that when \(T\) is high, the probability \(P\) is close to 1, allowing the algorithm to accept almost any move. This chaotic behavior enables the search to escape local minima—shallow valleys that would trap a greedy algorithm. As \(T\) decreases, \(P\) drops, and the algorithm becomes increasingly selective. Then, either \(T\) hits a cutoff, or our stopping criteria, \(M\), is met, and it settles into the lowest error state found.

\section{Handwriting Classification}

Handwriting classification with artificial intelligence starts by using data (a drawing of the letter `A') and labeling it (`A'). Then we use that labeled data to train an algorithm to recognize the letter `A' in drawings and classify them correctly. We can evaluate how good an algorithm is by giving it new drawings and asking it to classify them as a specific letter. In general, we would use a cross-validation set that partitions the data into training and evaluation sets with no overlap.

To classify data, we could try a specific algorithm, k-nearest neighbors, which uses a distance function to measure the similarity between a new drawing and our labeled data. Suppose that the drawings are composed on a \(40 \times 40\) drawing grid, where the pixels can be either on or off. Then, the model must determine which known examples are most similar to this new grid. This determination process can be performed by calculating the distance between the input value and each element in the training data. Then, these distance-label pairs are inserted into a priority queue that automatically orders them from nearest to farthest. Hence, we can access the most relevant data points without manually sorting the entire dataset. Once the algorithm organizes the distances, we can extract the top \(k\) elements from the priority queue, classify them as the nearest neighbors, and add them to a histogram. Now, we can take a vote to get the final classification by counting the frequency of each label among the \(k\) neighbors. The label with the highest count is the predicted classification. 

As we can see, the accuracy of data classification depends heavily on the value of the hyperparameter \(k\). If we set \(k\) too high, the model will misclassify a letter as another letter, but if we put it too low, the model will not generalize properly; this is known as overfitting. Thus, we are left with a hyperparameter optimization problem.

\subsection{Applying Optimizers}

As we alluded to in the previous section, we must optimize the hyperparameter \(k\) to maximize the histogram's voting accuracy. Since \(k\) is the only hyperparameter we are testing, grid search is a viable choice, as we can define a range for \(k\) (e.g., 1 to 20) and exhaustively check each integer in that range to find the highest accuracy. We would not want to try other hyperparameter optimizers, as we do not have enough hyperparameters to justify their use. If we had introduced another option, such as a different distance metric or threshold values for pixel transparency, the search space would grow, and something like random search would be viable. 

\section{Sentiment Analysis}

In natural language processing, we often want to understand someone's sentiment (how they feel) when they buy a product or talk to a friend. Our goal is to determine if they are positively or negatively affected by the product. With that information, we could recommend similar products through ads or suggest a different product to the customer, for example. To analyze sentiment, we use a similar process to handwriting classification: we gather data (e.g., Amazon reviews) and assign a positive or negative sentiment. Then we train the model on that labeled data and use it to classify new reviews.  There are many algorithms that can handle sentiment analysis; in fact, we could use the k-nearest neighbors algorithm again, but this process is generally done with a support vector machine, or SVM. 

If we want to classify data with an SVM, we must contend with multiple hyperparameters. The SVM has two sets of hyperparameters: kernel and regularization. The four most common kernels for SVM are linear, polynomial, radial basis function (RBF), and sigmoid. Then, the regularization values, \(c\) and \(\gamma\), are applied to the kernels in different ways. The hyperparameter \(c\) is only used for the polynomial and sigmoid kernels, whereas \(\gamma\) is used in every kernel except linear \parencite{jain_2024_svm}. Since we usually work with just two sentiments, it is common to pick the linear kernel over the others, but as mentioned in the introduction, if there are many hyperparameter possibilities, we must assume all of them are viable \parencite{reddy_2023_sentiment}. This raises another problem we can pursue with hyperparameter optimizers.

\subsection{Applying Optimizers}

For SVMs, we might first look to grid or random search. As previously noted, grid search can perform the best when dealing with categorical variables within a limited domain, such as our four kernel choices. With the addition of regularization parameters \(c\) and \(\gamma\), our search space increases, so we may be wary of grid search. However, we could leverage random search's low effective dimensionality property to address the space problem. With that being said, we could also try simulated annealing, as it may perform better in this space due to its ability to explore many combinations while still keeping track of the best performers. As mentioned earlier, simulated annealing works by choosing a new configuration that is close to the current one. Since we are using continuous parameters, this is relatively straightforward: we pick a value for \(c\) or \(\gamma\) that is slightly higher or lower than the current value. For the kernel type, we arbitrarily change to one of the three options. This flexibility allows us to explore discrete kernel choices while fine-tuning the continuous regularization parameters. Note that we could use Bayesian optimization here, but since simpler solutions would work, it is best to try them first. If we find that their results are poor, then we could consider a more sophisticated solution. 

\section{Image Classification}

In the previous section on handwriting classification, we greatly simplified the problem by restricting the domain to binary-valued letters. If we want to employ a generic image classifier that can handle any image, we must turn to convolutional neural networks (CNNs) to handle the increased detail. That is, we now have an enormous number of pictures, each containing arrays of pixel values for red, green, and blue, with possibly hundreds of thousands of unique labels.  

For this problem, we can leverage existing trained models, like VGG16, which contains 138 million parameters. This CNN was trained with the ImageNet database, which has over 14 million labeled training images \parencite{gabrielferrer_2025_csci}. We can configure this model so that we don't affect its trained parameters. Then, we can build another neural network at the end of VGG16 to perform specific image classification of our specific dataset (e.g., cats and dogs). However, we see that this addition introduces many hyperparameters. We must find an optimal rescaling factor to reduce the neural network's training time. In addition, we must decide how many layers to add and, for each layer, which activation function to use. In short, there are many hyperparameters we need to optimize, and that is not even accounting for training (e.g., learning rate, loss function, optimizer, etc.). 

\subsection{Applying Optimizers}

This is a prime problem for Bayesian optimization. We are handling a neural network that takes time to train, and (depending on available resources) is energy-expensive. We need the optimizer to be efficient and effective at picking the optimal hyperparameters. As discussed earlier, the optimizer treats the model's accuracy as a ``black box'' function that is expensive to evaluate. So, we will build a probabilistic model of how the layers and activation functions affect accuracy. In turn, the algorithm will learn from each run and improve incrementally. In this situation, every training run matters, so we cannot resort to exhaustive measures. 

\section{Conclusion}

Through this paper, we have explored four distinct hyperparameter optimizers: grid and random search, Bayesian optimization, and simulated annealing. We discussed that grid search is best when we have relatively few hyperparameters and spare resources for parallelization. Then, we looked at random search, which had many of the same properties as grid search but performed better in general. The better performance is a result of the random search's low effective dimensionality: since it samples hyperparameters uniformly, it covers a broader range of essential hyperparameters. This is because we showed that grid search wastes six tests for a \(3 \times 3\) grid. Then, we moved to Bayesian optimization, a complex algorithm that uses a probabilistic model of the objective function, employing prior distributions, past data, and acquisition functions. A specific acquisition function we looked into, upper-confidence bound (UCB), balances exploration and exploitation to make our search as effective as possible. We finished off with simulated annealing, an algorithm based on the physical process of annealing materials. This algorithm excels at avoiding minima and maintaining a broad search space. 

After we introduced each algorithm, we identified common problems in artificial intelligence and walked through each issue, discussing how each optimizer can be applied to address it. For the first problem, we looked into handwriting classification and found that, since we are using k-nearest neighbors, which has only one hyperparameter, grid search would be best. We could use grid search to exhaustively explore a range of possible \(k\) values, trying each one without worrying about the curse of dimensionality. Then, we moved on to sentiment analysis with support vector machines. We decided that grid and random search would not be adequate, as the variables are continuous. This left us with simulated annealing, as we showed that properties to explore the global space were the most satisfactory. This left us with the most complex problem: image classification. We went through an example of using VGG16, a robust convolutional neural network pretrained on millions of images. We wanted to build a neural network on top of this one, so we had to contend with at least six hyperparameters. This made Bayesian optimization the most appealing and the most cost-effective option. 

Therefore, we have shown through our exploration of hyperparameter optimizers and specific problems that as the complexity of training an algorithm increases, the need for advanced hyperparameter optimization methods also increases.

\printbibliography

\end{document}
